\section{Discussion}\label{sec:discussion}
Simulation-free variational inference promises the accuracy of simulation-based algorithms at a fraction of the cost, but only when the variational family is rich enough to represent the posterior.
We have shown that under high posterior uncertainty, existing simulation-free algorithms fall short of their promise: by committing to an arbitrary drift consistent with their one-time marginals, they restrict the posterior to a subset of the path laws available to simulation-based methods. 
Helmholtz-SDE removes this restriction. 
Starting from any reference drift, it adds a closed-form correction that preserves the prescribed marginals while optimizing over a larger family of path laws. 
The result is a simulation-free algorithm that closes the expressivity gap to simulation-based VI and yields more faithful recovery of the underlying dynamics.

Two limitations point to natural extensions. 
First, the analytical form for the Helmholtz correction depends on the state independence of the diffusion coefficient and Gaussianity of the marginals. 
Relaxing either assumption, for example by learning the divergence-free correction, would broaden the algorithm's applicability.
Second, higher-order polynomial corrections to the residual scale unfavorably as the latent dimension grows; in our experiments we did not observe gains beyond the linear and quadratic approximations, but more efficient parameterizations may be beneficial in settings we have not explored. 
More broadly, we view simulation-free VI as a promising path toward scalable unsupervised discovery of dynamical systems, and Helmholtz-SDE as a step toward making that path competitive with, and ultimately preferable to, its simulation-based counterpart.